\clearpage
\part{第三题：字符串解析与综合模拟}
\section{字符串解析与综合模拟基础}
\subsection{字符数组与字符串操作}
C风格字符串是以结束符\texttt{'\textbackslash0'}结尾的字符数组。长度不包括结束符，容量必须包括它。strlen求长度，strcmp比较内容，strcpy复制，头文件为\texttt{<cstring>}。
% snippet:c_string
\begin{lstlisting}
char a[20] = "cat", b[20];
int len = (int)strlen(a); // 3.
strcpy(b, a);            // Copies the final '\0' too.
b[0] = 'b';              // b is now "bat".
int relation = strcmp(a, b); // Positive: "cat" > "bat".
\end{lstlisting}
strcmp返回负数、0、正数，分别表示字典序小于、相等、大于；不能用数组名的==比较内容。复制前保证目标容量足够。strlen、strcpy随字符串长度线性增长；strcmp最坏需要检查完整公共前缀。循环处理字符时先保存长度，避免每轮重新调用strlen。

\subsection{整行输入与字符处理}
fgets最多读取容量减1个字符，保留读到的换行，并补结束符。若一行最多有1000个正文字符，下面的容量可以容纳正文、行尾和结束符。
% snippet:line_input
\begin{lstlisting}
char line[1003];
if (fgets(line, sizeof(line), stdin) != NULL)
    line[strcspn(line, "\r\n")] = '\0';
\end{lstlisting}
strcspn返回第一个回车或换行的位置，无行尾时返回字符串长度；该位置置0，只去掉行尾，不删除正文空格。若先scanf读数字，再fgets读下一行，应先消耗数字所在行的剩余字符：
\begin{lstlisting}
int c;
while ((c = getchar()) != '\n' && c != EOF) {}
\end{lstlisting}
字符分类函数位于\texttt{<cctype>}。isdigit判断数字，isspace判断空白，tolower转为小写；参数先转为unsigned char，避免有符号char的负值。
% snippet:char_class
\begin{lstlisting}
int length = (int)strlen(line), digits = 0;
for (int i = 0; i < length; i++) {
    unsigned char c = (unsigned char)line[i];
    if (isdigit(c)) digits++;
    line[i] = (char)tolower(c);
}
\end{lstlisting}
扫描长度L的文本为$O(L)$。若规则只把普通空格当分隔符，就比较\texttt{' '}；isspace还会接受制表符等空白。

\subsection{扫描指针与状态机}
扫描指针p表示下一个尚未处理的字符。连续数字应一次读取成一个数，而不是把每个字符单独当成数据。
% snippet:scan_number
\begin{lstlisting}
// s is a null-terminated string; p initially points at a digit.
long long value = 0;
while (s[p] >= '0' && s[p] <= '9') {
    value = value * 10 + s[p] - '0';
    p++;
}
\end{lstlisting}
例如从\texttt{12,7}的开头读取，value依次为1、12，结束时p指向逗号。数值范围必须允许每次乘10再加数字。

状态机用少量状态说明同一字符在不同上下文中的含义。例如引号外的空格是分隔符，引号内的空格是正文；用in\_quote记录当前是否在引号中。
% snippet:scan_state
\begin{lstlisting}
// Input has matched quotes and no escaping.
int in_quote = 0, separators = 0;
for (int p = 0; s[p] != '\0'; p++) {
    if (s[p] == '"') in_quote = !in_quote;
    else if (s[p] == ' ' && !in_quote) separators++;
}
\end{lstlisting}
\texttt{a "b c" d}的两个引号外空格被计数，引号内空格不计。每个字符触发一次状态转移，时间$O(L)$，额外空间$O(1)$。完整解析还需按规则决定空片段、转义和非法输入的处理方式。

\subsection{状态存储与更新顺序}
模拟用变量保存当前状态，按照操作规则更新。每个变量应对应明确的对象和时刻：编号作下标，多个属性放入结构体，查询读取该时刻已生效的状态。
\subsubsection{同步更新}
若一轮中的每个新值都依赖上一轮，必须从旧数组a读取，写入新数组b；全部算完后再复制回a。
% snippet:synchronous
\begin{lstlisting}
// One round: new value = old value + old immediate neighbors.
for (int i = 0; i < n; i++) {
    b[i] = a[i];
    if (i > 0) b[i] += a[i-1];
    if (i+1 < n) b[i] += a[i+1];
}
for (int i = 0; i < n; i++) a[i] = b[i];
\end{lstlisting}
$[1,2,3]$一轮后为$[3,6,5]$。若直接改a，计算中间位置时会读到已更新的3，结果变为8。每轮$O(n)$时间，两个数组$O(n)$空间。
\subsubsection{顺序更新}
若操作依次生效，后一条必须看到前一条的结果。事件模拟还要确定同一时刻的先后，例如先加入该时刻到达的对象，再从全部候选中选择。
\begin{lstlisting}
// Events are already sorted by arrival time.
while (next < n && events[next].arrival <= now) {
    add_candidate(events[next]);
    next++;
}
if (has_candidate()) execute_one();
\end{lstlisting}
小于等于把恰好在now到达的对象也纳入本次选择。每个事件只加入一次，加入阶段总计处理n项；选择和执行成本由相应数据结构决定。

\subsection{栈与递归解析}
\subsubsection{栈保存未完成状态}
处理嵌套结构时，最近打开的结构应最先关闭，用栈保存尚未配对的左括号。
% snippet:parse_stack
\begin{lstlisting}
// s contains only '(' and ')'; stack capacity >= strlen(s).
int top = 0, valid = 1;
for (int p = 0; s[p] != '\0'; p++) {
    if (s[p] == '(') st[top++] = p;
    else if (top == 0) { valid = 0; break; }
    else top--;
}
if (top != 0) valid = 0;
\end{lstlisting}
右括号遇到空栈表示没有对应左括号，扫描结束仍有元素表示左括号未关闭。栈中存位置还可以记录配对关系。时间$O(L)$，最坏空间$O(L)$。

\subsubsection{递归函数对应语法层级}
递归解析把语法中的一个结构交给一个函数。以非负整数、加号和括号为例：
\[
\begin{aligned}
\mathrm{expr}&=\mathrm{factor}\ \{+\ \mathrm{factor}\},\\
\mathrm{factor}&=\mathrm{number}\mid(\ \mathrm{expr}\ ).
\end{aligned}
\]
expr处理连续加法，factor处理数字或括号。共享指针p始终指向下一个未读字符；本片段假定语法合法、没有空白，数值运算不溢出。
% snippet:recursive_parser
\begin{lstlisting}
char s[1005]; int p = 0;
long long parse_expr();
long long parse_factor() {
    if (s[p] == '(') {
        p++;
        long long value = parse_expr();
        p++; // Consume the matching ')'.
        return value;
    }
    long long value = 0;
    while (s[p] >= '0' && s[p] <= '9')
        value = value*10 + s[p++] - '0';
    return value;
}
long long parse_expr() {
    long long value = parse_factor();
    while (s[p] == '+') {
        p++;
        value += parse_factor();
    }
    return value;
}
\end{lstlisting}
解析\texttt{1+(2+3)}时，外层读到左括号便递归求出5，再与1相加；内层在右括号处停止，由调用它的factor消费右括号。每个字符消费一次，时间$O(L)$，调用栈空间$O(H)$，H为嵌套深度。若题目允许非法输入，需要检查预期的字符，并在顶层确认整行读完。

\subsection{树形结构与依赖处理}
\subsubsection{父子关系与子树}
树形结构用父编号或孩子列表表达。普通数组可以用first[u]保存第一个孩子，next[v]保存同父节点下的下一个孩子。
% snippet:tree_build
\begin{lstlisting}
// Nodes 0..n-1; parent[root] = -1.
for (int u = 0; u < n; u++) first[u] = next[u] = -1;
for (int v = 0; v < n; v++) {
    int u = parent[v];
    if (u == -1) continue;
    next[v] = first[u];
    first[u] = v;
}
\end{lstlisting}
孩子按逆插入顺序连接；若题目要求输入顺序，需维护尾指针或按反向编号插入。递归遍历孩子后再处理父亲，可以汇总子树信息：
% snippet:tree_size
\begin{lstlisting}
int subtree_size(int u) {
    int total = 1;
    for (int v = first[u]; v != -1; v = next[v])
        total += subtree_size(v);
    return total;
}
\end{lstlisting}
返回值包含u本身。合法树上每个节点和父子边访问一次，时间$O(n)$，存储$O(n)$，调用栈深度为树高H；链状树可能达到n层。

\subsubsection{依赖顺序与环检测}
若u的结果依赖v，就应先完成v。依赖图不一定是树：多个对象可以依赖同一个对象，也可能出现环。用三种颜色记录进度：0未开始，1正在当前递归链中，2已经完成。
% snippet:dependency
\begin{lstlisting}
// dependency[u] lists objects needed by u; color is initially zero.
bool solve(int u) {
    if (color[u] == 1) return false;
    if (color[u] == 2) return true;
    color[u] = 1;
    for (int i = 0; i < (int)dependency[u].size(); i++)
        if (!solve(dependency[u][i])) return false;
    color[u] = 2;
    order[count++] = u;
    return true;
}
\end{lstlisting}
再次遇到颜色1，说明当前链回到尚未结束的对象，存在环；遇到颜色2则复用已经完成的结果，不是环。order中依赖对象先于使用它的对象，可以按此顺序计算。对所有未访问节点调用solve即可检查全图，n个对象、m条依赖的时间与存储为$O(n+m)$；实际结果计算和输出若更大，应另计成本。
